% Options for packages loaded elsewhere
\PassOptionsToPackage{unicode}{hyperref}
\PassOptionsToPackage{hyphens}{url}
\PassOptionsToPackage{dvipsnames,svgnames,x11names}{xcolor}
%
\documentclass[
  11pt,
]{article}
\usepackage{amsmath,amssymb}
\usepackage{iftex}
\ifPDFTeX
  \usepackage[T1]{fontenc}
  \usepackage[utf8]{inputenc}
  \usepackage{textcomp} % provide euro and other symbols
\else % if luatex or xetex
  \usepackage{unicode-math} % this also loads fontspec
  \defaultfontfeatures{Scale=MatchLowercase}
  \defaultfontfeatures[\rmfamily]{Ligatures=TeX,Scale=1}
\fi
\usepackage{lmodern}
\ifPDFTeX\else
  % xetex/luatex font selection
\fi
% Use upquote if available, for straight quotes in verbatim environments
\IfFileExists{upquote.sty}{\usepackage{upquote}}{}
\IfFileExists{microtype.sty}{% use microtype if available
  \usepackage[]{microtype}
  \UseMicrotypeSet[protrusion]{basicmath} % disable protrusion for tt fonts
}{}
\makeatletter
\@ifundefined{KOMAClassName}{% if non-KOMA class
  \IfFileExists{parskip.sty}{%
    \usepackage{parskip}
  }{% else
    \setlength{\parindent}{0pt}
    \setlength{\parskip}{6pt plus 2pt minus 1pt}}
}{% if KOMA class
  \KOMAoptions{parskip=half}}
\makeatother
\usepackage{xcolor}
\usepackage[margin=18mm]{geometry}
\usepackage{longtable,booktabs,array}
\usepackage{calc} % for calculating minipage widths
% Correct order of tables after \paragraph or \subparagraph
\usepackage{etoolbox}
\makeatletter
\patchcmd\longtable{\par}{\if@noskipsec\mbox{}\fi\par}{}{}
\makeatother
% Allow footnotes in longtable head/foot
\IfFileExists{footnotehyper.sty}{\usepackage{footnotehyper}}{\usepackage{footnote}}
\makesavenoteenv{longtable}
\setlength{\emergencystretch}{3em} % prevent overfull lines
\providecommand{\tightlist}{%
  \setlength{\itemsep}{0pt}\setlength{\parskip}{0pt}}
\setcounter{secnumdepth}{-\maxdimen} % remove section numbering
% Shared preamble for the pandoc-generated roadmap PDFs (pdflatex).
\usepackage[T1]{fontenc}
\usepackage[utf8]{inputenc}
\usepackage{lmodern}
\usepackage{microtype}
\usepackage{amssymb}
\usepackage{booktabs}
\usepackage{longtable}
\usepackage{array}
\usepackage{xcolor}

% Map the non-ASCII glyphs the markdown uses so pdflatex doesn't choke.
\usepackage{newunicodechar}
\newunicodechar{→}{$\rightarrow$}
\newunicodechar{←}{$\leftarrow$}
\newunicodechar{↔}{$\leftrightarrow$}
\newunicodechar{⇒}{$\Rightarrow$}
\newunicodechar{±}{$\pm$}
\newunicodechar{≈}{$\approx$}
\newunicodechar{≥}{$\geq$}
\newunicodechar{≤}{$\leq$}
\newunicodechar{×}{$\times$}
\newunicodechar{÷}{$\div$}
\newunicodechar{·}{$\cdot$}
\newunicodechar{…}{\dots}
\newunicodechar{²}{\textsuperscript{2}}
\newunicodechar{³}{\textsuperscript{3}}
\newunicodechar{°}{\textdegree}
\newunicodechar{µ}{$\mu$}
\newunicodechar{σ}{$\sigma$}
\newunicodechar{ε}{$\varepsilon$}
\newunicodechar{θ}{$\theta$}
\newunicodechar{τ}{$\tau$}
\newunicodechar{γ}{$\gamma$}
\newunicodechar{ω}{$\omega$}
\newunicodechar{ζ}{$\zeta$}
\newunicodechar{Σ}{$\Sigma$}
\newunicodechar{√}{$\surd$}
\newunicodechar{✓}{\checkmark}
\newunicodechar{•}{\textbullet}
\newunicodechar{≠}{$\neq$}
\newunicodechar{∞}{$\infty$}
\newunicodechar{ }{~}
% status emoji + checks -> short text (pdflatex has no emoji glyphs)
\newunicodechar{🔬}{{[study]}}
\newunicodechar{🔴}{{[planned]}}
\newunicodechar{🟡}{{[planned]}}
\newunicodechar{🟢}{{[easy]}}
\newunicodechar{🗺}{{[map]}}
\newunicodechar{✅}{{[done]}}
\newunicodechar{⚠}{{[!]}}
\newunicodechar{️}{}   % U+FE0F variation selector
\newunicodechar{🤖}{}
\ifLuaTeX
  \usepackage{selnolig}  % disable illegal ligatures
\fi
\IfFileExists{bookmark.sty}{\usepackage{bookmark}}{\usepackage{hyperref}}
\IfFileExists{xurl.sty}{\usepackage{xurl}}{} % add URL line breaks if available
\urlstyle{same}
\hypersetup{
  colorlinks=true,
  linkcolor={blue},
  filecolor={Maroon},
  citecolor={Blue},
  urlcolor={blue},
  pdfcreator={LaTeX via pandoc}}

\author{}
\date{}

\begin{document}

\hypertarget{study-live-drone-pid-autotuning-in-a-rig-driven-by-the-gcs}{%
\section{Study: live-drone PID autotuning in a rig, driven by the
GCS}\label{study-live-drone-pid-autotuning-in-a-rig-driven-by-the-gcs}}

Status: 🔬 study / feasibility. Next phase of the autotuner: reuse the
SITL tuning method to tune a \textbf{real} flight controller held in a
physical test rig, driven over the live GCS link. Builds on the autotune
stack (\texttt{navigator/src/autotune/}), the source state machine
(\url{gcs-source-state-machine.md}), and the
\href{../../../docs/telemetry/time_sync.md}{time-sync} work.

\hypertarget{context--goal}{%
\subsection{Context \& goal}\label{context--goal}}

Today the autotuner runs entirely in simulation: \texttt{AutotuneWorker}
spawns an isolated SITL (\texttt{vsim\_d} + \texttt{vayu\_sitl}), the
optimizer proposes gains, and a \texttt{runRollout()} excites the
simulated craft and scores the response. The operator then
\textbf{applies} the winning gains to a real board over the live link
(\texttt{CMD\_SET\_PID}).

The goal of this study: close the loop on \textbf{hardware} --- run the
same optimizer/cost/excitation against a \textbf{real drone in a rig},
scoring real telemetry, so the gains are tuned on the actual airframe +
motors + sensors rather than the sim model.

\hypertarget{feasibility-verdict}{%
\subsection{Feasibility verdict}\label{feasibility-verdict}}

\textbf{Feasible, and most of the stack is already reusable --- but it
needs one enabling firmware command and a deliberate safety design.} The
blocker is that the GCS has no way to inject stick/setpoint excitation
over the live link.

\hypertarget{reusable-as-is-backend-agnostic}{%
\subsubsection{Reusable as-is
(backend-agnostic)}\label{reusable-as-is-backend-agnostic}}

\begin{itemize}
\tightlist
\item
  \textbf{Optimizer} (\texttt{Optimizer.\{h,cpp\}}) and \textbf{Space}
  (\texttt{Space.\{h,cpp\}}) --- pure numeric search + parameter bounds.
  No sim knowledge.
\item
  \textbf{Cost} (\texttt{Cost.\{h,cpp\}}) --- pure function of a
  telemetry window (setpoint-vs-measured IAE + overshoot + chatter). The
  fields it needs are exactly the real FC\textquotesingle s
  \texttt{control\_telemetry\_t} (setpoints + measured + outputs),
  telemetered as \texttt{SYSTEM\_ORIGIN\_PID\_ERROR} and surfaced as
  \texttt{ControlLoopData}.
\item
  \textbf{AutotuneEngine} --- the rollout is already a
  \texttt{std::function\textless{}...\textgreater{}}
  (\texttt{AutotuneEngine.h:28}), i.e. \textbf{the backend is
  pluggable}. The engine has no SITL coupling.
\item
  \textbf{applyGains()} --- \texttt{setPid}/\texttt{setGyroLpf} are
  plain NavLink commands that work over the live link unchanged.
\end{itemize}

\hypertarget{the-seam-to-introduce}{%
\subsubsection{The seam to introduce}\label{the-seam-to-introduce}}

\texttt{runRollout()} is hard-typed to \texttt{SitlStack\ \&} (no
interface). To plug a live backend, abstract the \textasciitilde10 calls
it makes into an interface (or template):

\begin{verbatim}
setPid, setGyroLpf, setRc, arm, disarm, waitLevel, clearSamples, snapshot   // reusable
reset(seed), setTestRig(tetherK)                                            // SITL-only
\end{verbatim}

Introduce an \texttt{IExciteStack} interface implemented by both
\texttt{SitlStack} (today) and a new \textbf{\texttt{LiveRigStack}}
(drives the real FC over the GCS link). The SITL-only ops become no-ops
/ different strategies on the live backend (see "Sim vs real" below).

\hypertarget{the-enabling-gap-gcs-excitation-over-the-link}{%
\subsection{The enabling gap: GCS excitation over the
link}\label{the-enabling-gap-gcs-excitation-over-the-link}}

\textbf{There is no command to inject RC sticks or rate/angle setpoints
over the live link.} In SITL the tuner writes
\texttt{sim\_rc\_channels{[}{]}} via a PTY that replaces the iBus
receiver (\texttt{rc\_task.c}, \texttt{VAYU\_SIM} path); a real board
only takes RC from the physical receiver. Applying gains works
(\texttt{CMD\_SET\_PID}), but \textbf{exciting} the craft does not.

\textbf{Proposed firmware addition --- \texttt{CMD\_SET\_RC\_OVERRIDE}}
(a new GCS→FC command):

\begin{itemize}
\tightlist
\item
  Feeds GCS-supplied channel values into the same
  \texttt{ibus\_raw\_data} path the real receiver uses (so the entire RC
  → flight-mode → controller cascade is exercised, exactly like the SITL
  excitation). Tuning in \textbf{angle mode} then exercises the angle +
  rate loops the cost scores.
\item
  \textbf{Deadman/heartbeat:} the override is valid only for a short TTL
  (e.g. 200 ms) and must be refreshed; if GCS overrides stop (link drop,
  tuner aborts), the FC reverts to the physical receiver, which --- with
  no live stick --- trips the existing RC-loss failsafe → disarm. This
  makes "GCS goes away" fail safe by construction.
\item
  Must coexist with the physical receiver as a \textbf{safety override
  channel}: the operator\textquotesingle s TX arm switch / throttle and
  a hardware kill must always win.
\end{itemize}

This command is the one genuinely new, safety-critical piece. It
deserves its own focused design + review before any motors spin.

\hypertarget{rig-requirements}{%
\subsection{Rig requirements}\label{rig-requirements}}

The "test rig" in SITL is \textbf{pure sim physics}
(\texttt{physics\_core.cpp} hard-pin or soft critically-damped tether)
--- there is \textbf{no firmware rig mode}, and none is needed. The real
equivalent is \textbf{mechanical}:

\begin{itemize}
\tightlist
\item
  A \textbf{1-DOF or 3-DOF gimbal / tethered test stand} that frees
  rotation but constrains translation, so the craft can be excited in
  roll/pitch/(yaw) without flying away or needing altitude hold. The
  soft-tether sim model (\texttt{tetherK}) approximates a compliant rig;
  a stiff gimbal approximates the hard pin.
\item
  \textbf{Prop guards / containment} and a \textbf{physical kill} (cut
  power) independent of the firmware.
\item
  Tuning runs in \textbf{angle mode} so the firmware\textquotesingle s
  70° bank-angle failsafe (\texttt{angle\_controller.c},
  \texttt{MAX\_ANGLE\_CUTOFF}) catches a divergent excitation
  (it\textquotesingle s suppressed in acro).
\end{itemize}

\hypertarget{liverigstack--the-rollout-backend}{%
\subsection{LiveRigStack --- the rollout
backend}\label{liverigstack--the-rollout-backend}}

Maps each \texttt{IExciteStack} op onto the live link:

\begin{longtable}[]{@{}ll@{}}
\toprule\noalign{}
Op & Live implementation \\
\midrule\noalign{}
\endhead
\bottomrule\noalign{}
\endlastfoot
\texttt{setPid}/\texttt{setGyroLpf} &
\texttt{CommandCodec::encodeSetPid/GyroLpf} → \texttt{sendToFc} (queued
to the engine). \\
\texttt{setRc(...)} & \texttt{CMD\_SET\_RC\_OVERRIDE} frame at
\textasciitilde50 Hz with the deadman heartbeat. \\
\texttt{arm}/\texttt{disarm} & \texttt{CMD\_ARM}/\texttt{CMD\_DISARM}
(software-arm latch); confirm via system-state telemetry. \\
\texttt{snapshot}/\texttt{clearSamples} & Buffer the
\texttt{ControlLoopData} stream (the engine already decodes it); window
it per excitation. \\
\texttt{waitLevel} & Poll measured attitude from telemetry until
\textbar roll\textbar,\textbar pitch\textbar{} \textless{} ε (no instant
reset --- real settle). \\
\texttt{reset(seed)} & \textbf{No-op.} Replaced by "restabilize" ---
disarm/hold level between rollouts; no deterministic seed. \\
\texttt{setTestRig} & \textbf{No-op.} The mechanical rig is the
constraint. \\
\end{longtable}

\hypertarget{sim-vs-real--how-it-changes-the-algorithm}{%
\subsection{Sim vs real --- how it changes the
algorithm}\label{sim-vs-real--how-it-changes-the-algorithm}}

\begin{longtable}[]{@{}llll@{}}
\toprule\noalign{}
Property & SITL & Real rig & Consequence \\
\midrule\noalign{}
\endhead
\bottomrule\noalign{}
\endlastfoot
Reset & instant, deterministic (\texttt{reset(seed)}) & none --- must
restabilize & longer inter-rollout settle; no replay \\
Noise & seeded, reproducible & real, non-repeatable & rely on
\textbf{more repeats + statistics}; can\textquotesingle t average over
seeds \\
Constraint & sim tether & mechanical gimbal & rig compliance affects the
response (model the rig, or tune with it in-loop) \\
Iteration time & \textasciitilde seconds, fast & seconds--slower over
the link + settle & \textbf{smaller optimizer budget}; favor
sample-efficient search \\
Telemetry rate & high & \texttt{ControlLoopData} \textasciitilde18 Hz &
coarser cost window --- may need a higher tuning-telemetry rate during a
run \\
Timing & local & over WiFi/serial (jitter) & use the \textbf{time-sync}
offset to align the cost window to FC time; gate on link health \\
Failure & retry via reset & safety abort → auto-disarm & a
diverged/aborted rollout returns \texttt{nullopt}/\texttt{kBig}; the
loop must re-level before retry \\
\end{longtable}

\hypertarget{safety-architecture}{%
\subsection{Safety architecture}\label{safety-architecture}}

Firmware already provides strong backstops (verified):
\textbf{bank-angle failsafe (70°, angle mode)}, \textbf{RC-loss watchdog
(1 s) → FAILSAFE}, \textbf{throttle failsafe}, \textbf{motor cutoff on
any non-ARMED state}, and \textbf{CMD\_SET\_PID NaN/range validation +
persist}. The live-tuning design must add a GCS-side safety envelope on
top:

\begin{itemize}
\tightlist
\item
  \textbf{Deadman heartbeat} on \texttt{CMD\_SET\_RC\_OVERRIDE} (above)
  --- loss of GCS → revert to RX → RC-loss failsafe.
\item
  \textbf{GCS e-stop}: a single operator action that disarms
  immediately, plus an overall run timeout and a per-rollout watchdog.
\item
  \textbf{Throttle cap} during tuning (rig hover only); never command
  beyond a configured ceiling.
\item
  \textbf{Divergence → auto-disarm}: when the cost detects
  \textgreater limit angle error, stop exciting and disarm
  (don\textquotesingle t just score \texttt{kBig}).
\item
  \textbf{Operator-in-the-loop arming}: the physical TX arm switch and a
  hardware power kill always override the software latch.
\item
  A dedicated \textbf{safety review} of \texttt{CMD\_SET\_RC\_OVERRIDE}
  before first motor spin.
\end{itemize}

\hypertarget{state-machine-integration}{%
\subsection{State-machine integration}\label{state-machine-integration}}

Live tuning needs the GCS \textbf{in \texttt{Fc}} (live link, tx
allowed) \textbf{and} running a tuning loop --- unlike the current
\texttt{Autotune} state, which is an isolated SITL run that
doesn\textquotesingle t feed the engine. Options:

\begin{enumerate}
\def\labelenumi{\arabic{enumi}.}
\tightlist
\item
  \textbf{A new top-level state \texttt{RigTune}} --- entered from
  \texttt{Fc}, tx-allowed, feed active (it reads live telemetry to
  score). Strict-single-source teardown disarms + stops the override on
  exit. Cleanest fit for the FSM.
\item
  \textbf{A sub-mode within \texttt{Fc}} --- a background tuning task
  gated on \texttt{state()==Fc}. Less explicit, but no new state.
\end{enumerate}

Recommend (1): it makes the safety-critical mode explicit in the one
authority that already gates tx and the pill, and reuses the teardown
guarantees.

\hypertarget{design-decisions-round-2}{%
\subsection{Design decisions (round 2)}\label{design-decisions-round-2}}

Operator/safety model and run protocol, decided 2026-06-15:

\begin{itemize}
\tightlist
\item
  \textbf{Arming stays on the physical RC.} To enter \texttt{RigTune}
  the craft must already be \textbf{armed via the RC arm switch} with
  \textbf{both sticks held bottom-left} (safe idle). Only then can the
  GCS trigger the mode. Arm/disarm authority is never taken by the GCS.
\item
  \textbf{Decouple the iBus decode from the RC path.} Refactor the
  hardcoded iBus parse out of \texttt{rc\_task} into a clean
  \textbf{RC-source abstraction} so a GCS-injected source can be fed in
  without iBus frames "infecting" the GCS RC packets (and vice-versa).
  Safety vs excitation channels route separately (see new problems).
\item
  \textbf{Two exits:} the GCS can leave the mode, \textbf{or} the
  operator pulls \textbf{both sticks up-right} on the physical TX →
  FAILSAFE (hardware-side escape, independent of the GCS).
\item
  \textbf{Trim telemetry during a run:} disable all non-essential
  telemetry; only the cost-relevant stream flows, at a much higher rate.
\item
  \textbf{Pause clock discipline mid-eval:} stop time-sync corrections
  during an eval so the FC clock free-runs monotonically and the window
  is easy to serialize; the in-flight data carries a \textbf{GCS
  sub-millisecond timestamp}; \textbf{re-sync between runs} to stay
  wall-correct.
\item
  \textbf{Tune angle and rate modes separately} --- their gains can
  differ by mode.
\item
  More rollouts per eval for a usable approximation under real noise.
\end{itemize}

\hypertarget{additional-problems-found-in-review}{%
\subsection{Additional problems found in
review}\label{additional-problems-found-in-review}}

Beyond the round-2 decisions, these need answers before motors spin:

\begin{enumerate}
\def\labelenumi{\arabic{enumi}.}
\tightlist
\item
  \textbf{RC is per-purpose multiplexed, not all-or-nothing.} During the
  loop the \textbf{safety channels (arm switch, throttle, the up-right
  escape gesture) must keep coming from the physical RX}, while only the
  \textbf{excitation channels (roll/pitch/yaw, and the hover throttle)
  come from the GCS}. The RC-source refactor must define channel routing
  + precedence: \textbf{RX safety always vetoes} the GCS injection, and
  the escape gesture is evaluated on RX even while the GCS drives the
  excitation sticks.
\item
  \textbf{Throttle ownership.} Entering at bottom-left means min
  throttle (arm precondition met), but excitation needs hover thrust ---
  so the \textbf{GCS injects hover throttle} during the loop. The
  deadman must revert \emph{all} GCS channels to RX on loss, so throttle
  falls back to the operator\textquotesingle s min → safe settle.
\item
  \textbf{Excitation timing over a jittery link.} Streaming
  \textasciitilde50 Hz \texttt{CMD\_SET\_RC} is distorted by WiFi jitter
  (the chirp especially). \textbf{Strongly prefer FC-side parametric
  excitation}: the GCS sends a waveform descriptor (axis, step/chirp,
  amplitude, duration) and the FC plays it out at the loop rate,
  deterministically --- jitter-immune, low-bandwidth, reproducible.
  Bigger firmware change, but it removes the link from the
  timing-critical path.
\item
  \textbf{Cost alignment should use the FC\textquotesingle s
  already-paired (setpoint, measured).} \texttt{control\_telemetry\_t}
  carries \texttt{roll\_angle\_sp} \emph{and} \texttt{roll\_angle\_curr}
  sampled \textbf{together on the FC clock}. Score from those ---
  don\textquotesingle t try to align a GCS-side injected setpoint
  against FC-measured across the link, which the latency jitter would
  smear. Pausing discipline keeps that FC timestamp a stable monotonic
  reference within the eval; the GCS sub-ms stamp is for serialization +
  post-run wall mapping, not for the cost alignment itself.
\item
  \textbf{\texttt{CMD\_SET\_PID} persists to SD on every call.} A search
  applies gains hundreds of times → SD wear + write latency in the loop.
  Need a \textbf{volatile apply} (RAM-only) during tuning, persisting
  only the final accepted gains.
\item
  \textbf{Restore known-good gains on abort/exit.} A diverging rollout
  leaves the last (bad) gains live; on failsafe/exit the FC must fall
  back to a safe baseline, not whatever the optimizer last tried.
\item
  \textbf{Cascade order + per-mode storage.} The rate loop is inner,
  angle is outer --- \textbf{tune rate first, then angle} with rate
  fixed; tuning out of order is unstable. And "gains differ by mode"
  needs \textbf{per-flight-mode gain storage} in firmware (today
  \texttt{pid\_config} is \texttt{gains{[}ctrl{]}{[}axis{]}}, not
  per-mode) --- a prerequisite for separate angle/rate-mode tunes.
\item
  \textbf{Non-stationary plant.} Over hundreds of real rollouts the
  \textbf{battery sags and motors/ESCs heat up}, so the plant the
  optimizer sees drifts --- later rollouts aren\textquotesingle t
  comparable to earlier ones. Monitor pack voltage, cap session length /
  re-baseline, and prefer sample-efficient search.
\item
  \textbf{High-rate stream bandwidth + loss.}
  \texttt{control\_telemetry\_t} is 74 B; serial @230400 is
  \textasciitilde22.5 KiB/s, so a few-hundred-Hz stream is near the
  ceiling and UDP/WiFi can drop frames. Pick the max sustainable rate,
  prefer the \textbf{lossless serial} link for tuning, and
  \textbf{sequence-number + gap-detect} the samples --- reject any
  window with a gap rather than mis-scoring it.
\item
  \textbf{Optimizer robustness.} Real noise + drift make the cost much
  noisier than SITL; the current SPSA/structured search may converge
  poorly. Consider a \textbf{noise-aware / sample-efficient} optimizer
  (e.g. Bayesian optimization) for the live budget.
\item
  \textbf{Probing unstable gains is by design.} The search \emph{will}
  try bad gains → violent oscillation. Keep \textbf{conservative
  \texttt{Space} bounds}, a tight GCS-side divergence limit that
  \textbf{aborts + disarms fast} (ahead of the 70° firmware cutoff), and
  a soft-start.
\item
  \textbf{Re-level between rollouts (no reset).} Without
  \texttt{reset(seed)} the craft must settle on the gimbal between
  rollouts (recenter / brief level-hold or disarm) --- slower and
  occasionally needs operator help; budget for it.
\item
  \textbf{Rig fidelity is a fundamental limit.} Gimbal friction/inertia
  + the constraint bias the response, so rig-tuned gains may not be
  optimal in free flight (same caveat as the sim, different bias). Yaw
  especially may be constrained/low-authority on a tether.
\item
  \textbf{Sub-ms GCS clock = \texttt{QElapsedTimer}} (nanosecond
  monotonic) for intra-eval stamping, mapped to wall time at eval
  boundaries (\texttt{currentMSecsSinceEpoch} is ms-only).
\end{enumerate}

\hypertarget{refined-architecture-round-3--decided}{%
\subsection{Refined architecture (round 3 ---
decided)}\label{refined-architecture-round-3--decided}}

The review collapsed the design into a clean split, \textbf{GCS =
optimizer brain + operator UI + study/persistence; FC = per-pass
executor}:

\begin{itemize}
\tightlist
\item
  \textbf{Per-pass handshake (ACK\textquotesingle d, not streamed):} GCS
  sends the run plan (candidate gains + excitation descriptor) →
  operator physically stabilizes the craft on the rig and clicks
  \textbf{Start} in the GCS → FC applies the candidate \textbf{in RAM
  only (never SD)}, \textbf{generates the excitation itself at the loop
  rate}, captures the window, \textbf{computes the cost on the FC}, and
  sends back the cost + summary. GCS picks the next candidate. The
  inter-pass time naturally covers the round-trip while the operator
  re-levels.
\item
  \textbf{No high-rate link stream.} Normal telemetry keeps flowing for
  operator awareness; the high-granularity sensor + control-loop trace
  is \textbf{logged to SD on the FC} for offline study, recoverable
  later. This logging mode is explicit and must not disturb normal flow.
\item
  \textbf{Excitation is FC-internal}, overriding the target
  axis\textquotesingle s \emph{setpoint} during a pass; the
  \textbf{physical RC keeps arm + throttle + safety + the entry/exit
  gestures} (enter armed with both sticks bottom-left; exit by GCS, or
  both sticks up-right → FAILSAFE). The GCS does \textbf{not} stream RC
  --- so link jitter is off the timing-critical path entirely.
\item
  \textbf{Bank-angle cutoff disabled} during a pass (the craft is
  mechanically tied to the rig --- no runaway), scoped to the mode and
  auto-restored on exit. \texttt{Space} bounds still constrain
  candidates; the bounded pass duration + operator gesture are the
  abort.
\item
  \textbf{No intermediate persistence or commit.} Candidates run
  transiently; only an explicit GCS \textbf{Apply} (button +
  confirmation overlay) writes the chosen gains to SD and into the live
  loop. On exit the FC reverts RAM → committed gains.
\item
  \textbf{Clock:} pause time-sync discipline during a pass so the FC
  clock (used for the cost window + SD log) is clean + monotonic;
  re-sync between passes for wall-correct logs. FC-side cost means
  \textbf{no GCS sub-ms timestamp is needed}.
\item
  \textbf{Same SITL optimizer/cost}; angle and rate(acro) modes tuned as
  separate runs; optimizer selection is a later study if results
  disappoint.
\item
  \textbf{Operator-managed}: battery/thermal pause-or-continue is the
  operator\textquotesingle s call (not wired for automated decisions).
  Rig stability is the explicit goal; tethered (more-DOF, more-valid)
  tuning comes later, gated on rig success.
\end{itemize}

This resolves the round-2 problems \#1--\#4, \#8--\#14 by construction
(no RC streaming, no cross-clock alignment, no high-rate link, no
intermediate SD writes).

\hypertarget{final-clarifications-round-4--closes-the-open-questions}{%
\subsection{Final clarifications (round 4 --- closes the open
questions)}\label{final-clarifications-round-4--closes-the-open-questions}}

\begin{itemize}
\tightlist
\item
  \textbf{Sequencing:} this is a \emph{study only}. Execution starts
  \textbf{after NavLink v2 integration lands}; any new packets/features
  the autotuner needs are added to NavLink v2 \textbf{before} autotune
  execution begins. The reliable command/result exchange is therefore
  not a blocker to design now --- it\textquotesingle s the first
  execution step.
\item
  \textbf{Cost lives only on the FC.} The GCS never computes cost; it
  only helps the operator \emph{gauge} the run (plots/progress). No
  shared/duplicated cost to keep in sync.
\item
  \textbf{Statistical acceptance:} run \textbf{multiple rollouts per
  candidate and average} the response of the \emph{same} PID set until a
  \textbf{rig-stability criterion} is met --- this absorbs real-world
  noise (mirrors the SITL \texttt{repeats}).
\item
  \textbf{User-seeded start:} the operator provides the \textbf{initial
  gain seed} (per mode), so the search starts from a sane point rather
  than probing wild values.
\item
  \textbf{Phased \& pragmatic:} work with what exists;
  don\textquotesingle t pre-build abstractions. The firmware pieces
  below are sequenced into the execution phases, not all up front.
\end{itemize}

\hypertarget{still-open--firmware-deliverables-execution-phase-post-navlink-v2}{%
\subsection{Still open --- firmware deliverables (execution phase,
post-NavLink-v2)}\label{still-open--firmware-deliverables-execution-phase-post-navlink-v2}}

The architecture moves real work onto the FC; these are the new pieces,
all sequenced after NavLink v2 --- none is a study blocker:

\begin{enumerate}
\def\labelenumi{\arabic{enumi}.}
\tightlist
\item
  \textbf{Reliable plan/result/start exchange.} NavLink today is
  fire-and-forget. The tuning command + result handshake needs reliable
  (seq/ack/retry) delivery --- \textbf{build it on the in-progress
  NavLink v2 codec}, not a parallel layer. This is \textbf{gated on the
  NavLink v2 work} (\texttt{vayu-navlink-v2} worktree): RX/uplink demux
  (their Phase 3) and the COMMAND migration (Phase 4) are prerequisites
  for new reliable uplink commands. Coordinate the rig-tune messages
  there.
\item
  \textbf{FC pass-executor subsystem} (new, explicit/gated firmware
  task): receive plan → transient RAM gain apply (atomic, no
  half-applied set) → wait for Start → run excitation → window + cost →
  SD log → report → revert-to-committed on exit. Must coexist with
  flight code without disturbing it.
\item
  \textbf{Port the cost + excitation generator to firmware} (step
  doublet / chirp at loop rate; \texttt{axisCost}/\texttt{yawRateCost}).
  Share one implementation with the GCS/SITL path so live and sim costs
  stay identical --- don\textquotesingle t fork them.
\item
  \textbf{Decoupled RC path} (your point): factor the iBus decode behind
  an RC-source abstraction with a clean \textbf{setpoint-override hook}
  for the excitation, so iBus frames and the injected excitation
  can\textquotesingle t cross-contaminate; plus the entry preconditions
  gate (armed + sticks safe) and the scoped cutoff-disable that
  auto-reverts.
\item
  \textbf{Buffered async SD logging} at control-loop granularity that
  never stalls the loop; SD throughput bounds the log rate.
\item
  \textbf{FC real-time budget} check: control loop + excitation + cost
  accrual + SD log + normal telemetry within one pass must fit the
  schedule.
\end{enumerate}

\hypertarget{phased-plan}{%
\subsection{Phased plan}\label{phased-plan}}

\begin{enumerate}
\def\labelenumi{\arabic{enumi}.}
\tightlist
\item
  \textbf{Refactor to a backend interface} (no behavior change): extract
  \texttt{IExciteStack}; \texttt{SitlStack} implements it;
  \texttt{runRollout()} takes the interface. Unit-testable with a mock
  stack.
\item
  \textbf{\texttt{CMD\_SET\_RC\_OVERRIDE} firmware command} + deadman,
  behind a build/config flag. Bench-test on the board (no props): verify
  override drives the RC path, heartbeat-loss reverts to RX, arm switch
  overrides.
\item
  \textbf{\texttt{LiveRigStack}} over the link; dry-run the rollout
  \textbf{disarmed} (apply gains, send excitation, capture telemetry,
  compute cost) --- no motors.
\item
  \textbf{GCS safety envelope} (e-stop, timeouts, throttle cap,
  divergence-disarm) + the \texttt{RigTune} FSM state.
\item
  \textbf{Props-on rig bring-up}: single-axis (roll only), conservative
  budget, operator hand on the kill. Compare converged gains to the
  sim\textquotesingle s.
\item
  \textbf{Full} roll/pitch/(yaw), tune the budget/repeats for real-noise
  statistics.
\end{enumerate}

\hypertarget{open-questions--risks}{%
\subsection{Open questions / risks}\label{open-questions--risks}}

\begin{itemize}
\tightlist
\item
  \textbf{\texttt{CMD\_SET\_RC\_OVERRIDE} is safety-critical} --- it
  lets software move the sticks of an armed drone. Needs its own design
  + review; the deadman + RX-override + hardware kill are
  non-negotiable.
\item
  \textbf{Rig fidelity}: a stiff gimbal hides thrust-tilt (sim notes
  this over-tunes); a compliant tether is more honest but adds rig
  dynamics into the response.
\item
  \textbf{Telemetry rate / timing}: 18 Hz \texttt{ControlLoopData} may
  be too coarse for the cost window; consider a higher tuning-telemetry
  rate + time-sync alignment.
\item
  \textbf{Non-determinism}: real noise means more repeats and a
  statistical acceptance test ("is B really better than A?") rather than
  a single rollout.
\item
  \textbf{Optimizer budget}: real rollouts are slow; favor
  sample-efficient methods.
\end{itemize}

\hypertarget{key-references}{%
\subsection{Key references}\label{key-references}}

\begin{itemize}
\tightlist
\item
  Autotune:
  \texttt{navigator/src/autotune/\{AutotuneEngine,Optimizer,Space,Cost,Rollout,SitlStack,AutotuneWorker\}.\{h,cpp\}}
\item
  Excitation/RC path: \texttt{Rollout.cpp}
  (\texttt{exciteAxis}/\texttt{exciteOnce}), \texttt{SitlStack.cpp}
  (\texttt{setRc}/\texttt{rcWriterLoop}), firmware
  \texttt{src/comm/rc\_task.c} (\texttt{VAYU\_SIM} override),
  \texttt{src/comm/rc\_safety.c}
\item
  Commands: \texttt{navigator/src/protocol/CommandCodec.\{h,cpp\}},
  firmware \texttt{include/comm/comm\_types.h},
  \texttt{src/comm/comm\_processor.c},
  \texttt{src/control/pid\_config.c}
\item
  Telemetry for cost: \texttt{control\_telemetry\_t}
  (\texttt{include/variables.h}) → \texttt{ControlLoopData}
  (\texttt{navigator/src/core/Types.h})
\item
  Sim rig physics: \texttt{sim/vsim/src/physics\_core.cpp},
  \texttt{sim/vsim/include/vsim\_proto.h}
  (\texttt{vsim\_ctl\_testrig\_t})
\item
  FSM: \texttt{navigator/src/core/SourceController.\{h,cpp\}},
  \texttt{SourceState.h}
\end{itemize}

\hypertarget{changelog}{%
\subsection{Changelog}\label{changelog}}

\begin{longtable}[]{@{}lll@{}}
\toprule\noalign{}
Date & Author & Description \\
\midrule\noalign{}
\endhead
\bottomrule\noalign{}
\endlastfoot
15/06/2026 & ragnar-vallhala & Initial feasibility study \\
\end{longtable}

\end{document}
